\documentclass[11pt,a4paper]{article}
\usepackage[utf8]{inputenc}
\usepackage[italian]{babel}
\usepackage{geometry}
\usepackage{titlesec}
\usepackage{enumitem}
\usepackage{booktabs}
\usepackage{tabularx}
\usepackage{color}
\usepackage{colortbl}
\usepackage{tcolorbox}
\usepackage{fancyhdr}
\usepackage{graphicx}

\geometry{margin=1.5cm}

\definecolor{fedblue}{RGB}{15, 45, 94}
\definecolor{fedlightblue}{RGB}{26, 62, 120}
\definecolor{fedgold}{RGB}{255, 204, 0}

\titleformat{\section}
  {\normalfont\Large\bfseries\color{fedblue}}
  {}{0em}{}[\titlerule]

\titleformat{\subsection}
  {\normalfont\large\bfseries\color{fedblue}}
  {}{0em}{}

\pagestyle{fancy}
\fancyhf{}
\renewcommand{\headrulewidth}{1pt}
\renewcommand{\footrulewidth}{1pt}
\fancyhead[L]{USS Venture}
\fancyhead[C]{Star Trek Adventures}
\fancyhead[R]{Federazione}
\fancyfoot[C]{\thepage}

\begin{document}

\begin{center}
\Huge\textbf{\textcolor{fedblue}{USS Venture (Classe Nova)}}\\[0.5cm]
\Large\textbf{Pannello Master Star Trek Adventures}
\end{center}

\begin{tcolorbox}[colback=fedlightblue!20, colframe=fedblue, title=\textbf{Informazioni sulla Nave}]
\begin{tabular}{ll}
\textbf{Classe}: & Nova \\
\textbf{Equipaggio}: & 80-85 membri \\
\textbf{Velocità massima}: & Curvatura 8 \\
\textbf{Armamento}: & Phaser tipo X, tubi lanciasiluri \\
\textbf{Specializzazione}: & Ricerca scientifica e esplorazione \\
\end{tabular}
\end{tcolorbox}

\section{Quantum Data Points (QDP)}

\begin{tcolorbox}[colback=fedlightblue!10, colframe=fedblue]
\begin{tabularx}{\textwidth}{|X|X|X|}
\hline
\rowcolor{fedblue!20} \textbf{QDP Alpha} & \textbf{QDP Beta} & \textbf{QDP Gamma} \\
\hline
\large $\square\square$ \textbf{2} & \large $\square\square$ \textbf{0} & \large $\square\square$ \textbf{0} \\
\hline
\textbf{Conoscenza, Analisi,} & \textbf{Diplomazia, Comunicazione,} & \textbf{Risorse, Materiali,} \\
\textbf{Scienza} & \textbf{Cooperazione} & \textbf{Estrazione} \\
\hline
\end{tabularx}
\end{tcolorbox}

\section{Filosofia di Gioco}

\begin{tcolorbox}[colback=fedlightblue!10, colframe=fedblue]
\begin{itemize}[leftmargin=*]
\item \textbf{Esplorazione, comprensione, protezione} come valori fondamentali
\item Bilanciamento tra curiosità scientifica e cautela
\item Risoluzione dei conflitti preferibilmente non-violenta
\item Applicazione dei principi della Federazione: non-interferenza quando possibile
\end{itemize}
\end{tcolorbox}

\section{Meccaniche Chiave di Star Trek Adventures}

\begin{tcolorbox}[colback=fedlightblue!10, colframe=fedblue]
\begin{itemize}[leftmargin=*]
\item \textbf{Tiri Base}: Attributi + Discipline contro Difficoltà (tipicamente 2-3)
\item \textbf{Momentum}: Generato da successi eccezionali, spendibile per effetti positivi
\item \textbf{Determinazione}: Spesa invocando Valori per azioni decisive
\item \textbf{Assistenza}: Bonus di +1 dado quando membri dell'equipaggio collaborano
\item \textbf{Complicazioni}: Possono verificarsi su fallimenti critici (20 sui dadi)
\end{itemize}
\end{tcolorbox}

\section{Esempi di Difficoltà per Tiri Comuni}

\begin{tabularx}{\textwidth}{|l|X|c|}
\hline
\rowcolor{fedblue!20} \textbf{Azione} & \textbf{Attributo + Disciplina} & \textbf{Difficoltà} \\
\hline
Scansioni standard & Ragione + Scienza & 1 \\
\hline
Analisi tecnologia iconiana & Ragione + Scienza & 3 \\
\hline
Comunicazioni diplomatiche & Presenza + Comando & 2 \\
\hline
Modifiche ai sensori & Controllo + Ingegneria & 2 \\
\hline
Percepire intenzioni aliene & Intuito + Connessione & 3 \\
\hline
\end{tabularx}

\section{Fasi dell'Avventura}

\begin{tcolorbox}[colback=fedlightblue!10, colframe=fedblue]
\begin{enumerate}[leftmargin=*]
\item \textbf{Introduzione (15 minuti)}: Stabilire situazione e personaggi
\item \textbf{Missioni Iniziali (60 minuti)}: "Primo Contatto Iconiano" - Scansione dell'anomalia
\item \textbf{Svolta (90 minuti)}: "Risveglio del Portale" - L'anomalia si intensifica
\item \textbf{Conclusione (75 minuti)}: "La Scelta Finale" - Decidere il destino dell'artefatto
\end{enumerate}
\end{tcolorbox}

\section{Enigmi della USS Venture}

\begin{tcolorbox}[colback=fedlightblue!10, colframe=fedblue]
\begin{description}[leftmargin=*]
\item[Il Codice Biologico] I membri telepatici dell'equipaggio percepiscono una sequenza ripetitiva di emozioni che corrispondono a un codice iconiano.
\item[Sincronizzazione Quantica dei Sensori] L'artefatto esiste in tre stati quantici simultaneamente. Richiede tre ufficiali che lavorano in sincronia.
\item[La Biblioteca Vivente] Il simbionte della Dottoressa Krell inizia a ricordare dettagli di un incontro con gli Iconiani avvenuto in una vita precedente.
\item[Il Teorema Subspaziale] I computer risolvono spontaneamente complesse equazioni subspaziali. Richiede un approccio transdisciplinare.
\end{description}
\end{tcolorbox}

\section{Personaggi Chiave}

\begin{tcolorbox}[colback=fedlightblue!10, colframe=fedblue]
\begin{description}[leftmargin=*]
\item[Capitano Elara T'Nae] Umana/Vulcaniana, ex archeologa xenobiologa
\item[Comandante Sovak] Primo Ufficiale vulcaniano, logico ma preoccupato per i Romulani
\item[Tenente Comandante Reese] Capo Ingegnere, prodigio tecnologico
\item[Dottoressa Krell] Ufficiale Scientifico trill con simbionte antico
\item[Vorin] Ufficiale delle comunicazioni betazoide con abilità telepatiche
\end{description}
\end{tcolorbox}

\section{Principi della Federazione da Considerare}

\begin{tcolorbox}[colback=fedlightblue!10, colframe=fedblue]
\begin{itemize}[leftmargin=*]
\item Non-interferenza vs. responsabilità galattica
\item Condivisione di conoscenza vs. sicurezza
\item Cooperazione multiculturale vs. cautela strategica
\item Preservazione di culture aliene e manufatti storici
\item Priorità alla vita e alla sicurezza sopra l'avanzamento scientifico
\end{itemize}
\end{tcolorbox}

\section{Opzioni per la Fase Finale}

\begin{tcolorbox}[colback=fedlightblue!10, colframe=fedblue]
\begin{description}[leftmargin=*]
\item[Contenimento Scientifico] (5 QDP Alpha + 3 QDP Beta) Utilizzare campi di forza per stabilizzare l'artefatto per studi futuri
\item[Protocollo di Conservazione] (7 QDP Alpha) Spostare l'artefatto fuori fase per preservarlo in modo sicuro
\item[Integrazione Controllata] (4 QDP Alpha + 4 QDP Beta) Lavorare con altre fazioni per un protocollo di accesso condiviso
\end{description}
\end{tcolorbox}

\section{Meccanica del Teletrasporto in Presenza di Interferenze Iconiane}

\begin{tcolorbox}[colback=fedlightblue!10, colframe=fedblue]
\begin{itemize}[leftmargin=*]
\item \textbf{Costo Base}: 3 QDP 
\item \textbf{Capacità Unica}: "Biofiltro Avanzato" - Può rilevare tentativi di infiltrazione (2 QDP)
\item \textbf{Bonus}: +2 alla stabilità se si spendono QDP Alpha (esperienza con teletrasporti)
\item \textbf{Risoluzione}: 2d6 + modificatori, dove 10+ è successo completo
\end{itemize}
\end{tcolorbox}

\section{Promemoria e Suggerimenti Narrativi}

\begin{tcolorbox}[colback=fedlightblue!10, colframe=fedblue]
\begin{itemize}[leftmargin=*]
\item "I vostri tricorder rilevano fluttuazioni temporali attorno all'artefatto che le altre navi non sembrano notare..."
\item "L'ufficiale scientifico nota similarità tra questi simboli iconiani e frammenti studiati dall'Accademia delle Scienze..."
\item "Ricevete una comunicazione criptata da Narendra Station con nuove direttive riguardo l'artefatto..."
\item Concludere con discorsi ispirati "alla Picard" nei momenti cruciali
\item Enfatizzare l'ottimismo di Star Trek anche nelle situazioni difficili
\end{itemize}
\end{tcolorbox}

\section{Comandi del Bot}

\begin{tcolorbox}[colback=fedlightblue!10, colframe=fedblue]
\begin{description}[leftmargin=*]
\item[/register] Registrazione come Master
\item[/status] Stato corrente dell'avventura
\item[/scan] Eseguire una scansione dell'artefatto
\item[/qdp\_status] Vedere lo stato dei tuoi QDP
\item[/add\_qdp <tipo> <quantità> <descrizione>] Aggiungere QDP (alpha, beta, gamma)
\item[/spend\_qdp <tipo> <quantità> <descrizione>] Spendere QDP
\item[/resolve\_enigma\_piece <tipo>] Risolvere parte di un enigma
\item[/teleport <tavolo> <descrizione>] Tentare un teletrasporto
\item[/convergence\_roll <descrizione>] Eseguire un tiro di Convergenza nella fase finale
\item[/send\_message <tavolo> <messaggio>] Inviare un messaggio ad un altro tavolo
\end{description}
\end{tcolorbox}

\end{document}